\paragraph{The complete finite matrix spectrum and its exact spatial algebra.}
The 78-line local down-addition result extends to all eight source columns
$V_x=c_{x\downarrow}^\dagger G$.  On the entire 560-state invariant
sector, exact arithmetic proves $q(16H)V=0$ for the same monic degree-78
polynomial $q$ and verifies all 78 matrix moments in
$\operatorname{span}\{I,A\}$, where $A=A^\mathsf{T}$,
$A^2=3I$ and $\operatorname{tr}A=0$.
Polynomial division of the resolvent then gives
\[
\mathcal G(z)=
 \frac{16\bigl[p_I(x)I+p_A(x)A\bigr]}{1296q(x)},
\qquad x=16z-4\sqrt3.
\]
The numerator coefficients are exact convolutions of $q$ and those
source moments.  This is a rational resolvent identity, rather than a
truncated large-frequency expansion.

A second exact polynomial calculation gives
$q\mid(p_I^2-3p_A^2)$.  At each simple root $x_i$, define
$\alpha_i=p_I(x_i)/(1296q'(x_i))$ and
$\beta_i=p_A(x_i)/(1296q'(x_i))$.
The earlier interval certificates give $\alpha_i>0$; the new identity
gives $3\beta_i^2=\alpha_i^2$, and strict interval signs determine
$s_i=\operatorname{sign}\beta_i$.  Therefore
\[
R_i=\alpha_i I+\beta_i A=2\alpha_i P_{s_i},
\qquad P_\pm=\tfrac12(I\pm A/\sqrt3).
\]
Both $P_\pm$ have rank four.  There are 40 plus-band lines and 38
minus-band lines; each band's $2\alpha_i$ weights sum to one.
The rank of $E_iV$ equals that of its Gram matrix $R_i$, so the reducing
cyclic space generated by all eight probes has dimension
$78\cdot4=312$.  The remaining 248 source dimensions are orthogonal to
these probes and their time translates; their spectrum is not claimed.

Writing $\mathcal G=FI+JA$, both band functions have strictly negative
imaginary part for $\operatorname{Im}z>0$.  Hence
\[
\mathcal G^{-1}=\frac{FI-JA}{F^2-3J^2},
\qquad
\Sigma=zI-(A/16+I/2)-\mathcal G^{-1}.
\]
The full finite inverse response and memory are exactly onsite plus the
original signed nearest-neighbor adjacency.  This special cube algebra
does not supply an infinite-volume spatial limit or a Lorentz metric.

The certified common spectral support lies strictly between
$g=299/10000$ and $b=37/20$.  On the reducing cyclic space, both
$K=H-E_G$ and its compression perpendicular to the observed columns obey
these bounds.  The exact coupling $W=KV-VM_1$ lies in that same space,
so the orthogonal unobserved sector contributes zero to the Schur memory.
Using $W^\dagger W=I/4$, its order-$p$ Taylor remainder satisfies
\[
\|R_p\|\leq\frac{r^{p+1}}{4g^{p+1}(g-r)},
\qquad |z|\leq r<g.
\]
This improves the earlier conservative full-sector disk without
assigning the stronger gap to the unresolved sector or to infinite volume.
Evaluation at $x=-4\sqrt3$ now computes $-\mathcal G(0)$ and
$\Sigma(0)$ exactly over $\mathbb Q(\sqrt3)$, including the previously
bounded coefficients $12.9058381883$ and $-4.9026412320$.
No floating linear solve, native detector execution, physical species
assignment or relativistic rest-mass inference is used in this extension.

\paragraph{Excluding a complete finite-gap mass dictionary.}
The full support gives an exact test of assigning these visible
single-charge gaps to charged-lepton rest energies under one positive
unit factor.  Its largest-to-smallest ratio is below 62, including the
positive up-channel line.  The retained retrospective PDG 2025 benchmark
has muon/electron ratio above 206 and tau/electron ratio above 3477,
even in the declared central-value plus/minus reported-uncertainty boxes.
If a source ratio is at most $R$ and the target ratio at least $Q>R$,
simultaneous relative errors at most $\epsilon<1$ require
\[
 \frac{1-\epsilon}{1+\epsilon}\leq\frac RQ,
 \qquad \epsilon\geq\frac{Q-R}{Q+R}.
\]
The exact rational calculation therefore forces an error above 54 percent
for at least one member of an electron/muon assignment, and above
96 percent for an electron/tau assignment, for every common positive
scale and every choice of visible lines.  The target values do not enter
the source law or spectrum.  This rejects the stated finite-gap mass
dictionary, not a continuum, an unobserved or composite sector, or OPH
as a whole.  The finite gaps are not thereby proved to be relativistic
rest masses, and no eligible prospective comparison is added.
